\chapter{M12: konservative Schnittstelle und einseitige Mechanik}\label{api:fsi}
Das experimentelle Paket \texttt{fem\_toolbox.coupling} überträgt Lasten
zwischen extrudierten 2D-Fluidsegmenten und äußeren Triangle3-Facetten
einer Tet4-Struktur. Breite und Normalen sind explizit; strukturelle
Normalen werden aus dem eindeutigen Volumeneigentümer gewonnen.
Dokument 90 nennt Formen, Einheiten, Fehlerpfade und Reproduktion.

Die Überdeckung entsteht durch konvexen Polygonclip. Eine symmetrische
Dreipunktregel integriert auf jedem Teildreieck den Grad zwei exakt.
Baryzentrische P1-Gewichte bilden die skalare Matrix \(W\); für Vektoren
gilt \(H=W\otimes I_3\). Damit gehören Verschiebung und integrierte Kraft
zum selben überprüfbaren Paar:
\[
 u_q=H u_s,\qquad F_s=H^T F_q,\qquad F_q=t_q A_q.
\]
Konstantenreproduktion erhält die Kraft, Koordinatenreproduktion das Moment;
die Transposition erhält virtuelle Arbeit. Geprüfte Rohfehler haben Einheiten
N, N m und J. Relatives Gate ist \(10^{-10}\), absolute Floors sind jeweils
\(10^{-12}\) in der zugehörigen SI-Einheit.

\texttt{build\_domain\_interface} verlangt eine benannte FSI-Wand und übernimmt
ihre Breite aus der FluidDomain. \texttt{build\_interface\_map} nimmt die
Geometrie direkt entgegen. Vollständige beidseitige Überdeckung, entgegengesetzte
Normalen und Koordinatenreproduktion sind Pflicht. Lücken und Doppelungen
werden auch dann abgelehnt, wenn ihre Gesamtflächen sich aufheben.
Geometrische Toleranz ist \(10^{-10}L_{ref}\), relatives Flächengate
\(10^{-10}\). Es gibt keine Extrapolation oder Nachbarpunktverteilung.

Die Lastschnittstelle erwartet \(\sigma_f n_f\) in Pa und bildet
\(t_s=-\sigma_f n_f-p_{external}n_f\). Fluid- und äußerer Druckanteil
werden getrennt berichtet. Ein Druckdatumwechsel verlangt denselben Wechsel
des äußeren Drucks. Ergebnisarrays sind unveränderlich; eine herausgegebene
Matrix ist eine Kopie. Allgemeine virtuelle 3D-Wege dienen der Arbeitsprüfung.
Physikalische Netzbewegung verlangt dagegen eine ebene spannweiteninvariante,
von den Fluidsegmenten aufgelöste P1-Spur.

\texttt{solve\_one\_way} übergibt konservative Knotenkräfte an den bestehenden
statischen Struktursolver, einschließlich vorhandener Modelllasten und
Randwerte. Grenzen sind \(\max|u|/L_{ref}\leq0.02\) und
\(\max\|\nabla u\|_F\leq0.02\). Analytische konstante und affine Flächenlasten
prüfen alle Kräfte und Momente auf nicht passenden Netzen. Ein echter
Fluiddrucklauf liefert beim geführten elastischen Block die analytische
Verschiebung und ausgeglichene Lagerreaktionen. Dies ist noch kein
bidirektionaler FSI-Nachweis; Netzbewegung und Kopplungsiteration folgen separat.
